\chapter{Introduction}
\label{ch:introduction}

\section{Le problème de départ}

Quand une administration française met en production une
application qui traite des données sensibles, elle doit la faire
\emph{homologuer}. Concrètement, une autorité dans l'organisation
signe un document qui dit~: \emph{«~j'ai compris les risques, voici
les mesures que j'ai mises en place, et j'accepte ce qu'il reste~».}
Ce document s'appelle le \textbf{dossier d'homologation}. Le cadre
est posé par le RGS, et l'ANSSI publie un guide en neuf étapes pour
le construire.

Le souci, c'est qu'une fois le dossier signé, il est figé. Pendant
ce temps, le SI continue à vivre~:

\begin{itemize}
  \item de nouvelles CVE sortent chaque semaine sur les composants
        utilisés~;
  \item les règles réseau (firewalls, NSG Azure) sont modifiées
        pour les besoins du quotidien~;
  \item l'Active Directory bouge en permanence (arrivées,
        départs, changements de groupe)~;
  \item les règles de détection du SIEM sont ajustées, désactivées,
        oubliées.
\end{itemize}

Au bout de quelques mois, l'écart entre ce que dit le dossier et
ce que fait vraiment le SI devient gênant. À un an, on a vu des
audits où plus de 50\,\% des règles de filtrage ne correspondaient
plus à ce qui était documenté. C'est ce qu'on appelle dans la
suite l'\textbf{audit drift}.

\section{Ce qu'on fait dans ce mémoire}

L'idée du PFE est simple~: si la majorité de l'information qui
peuple un dossier d'homologation peut être collectée par script
(état des règles NSG, liste des CVE des images Docker, comptes AD,
règles Wazuh actives), alors on peut maintenir une version
\emph{tenue à jour} du dossier en automatisant cette collecte.

Pour le démontrer, le projet comporte deux livrables~:

\begin{description}
  \item[Un lab d'expérimentation, UrbanUpC.] Un SOC complet déployé
    sur Azure, avec un Active Directory, des applications web
    vulnérables, et un SIEM Wazuh qui détecte les attaques. Le
    lab sert à la fois de cible d'étude (\emph{«~de quoi parle
    un dossier d'homologation, concrètement~?~»}) et de banc de
    test pour l'outil.
  \item[Un outil d'automatisation, HOMO-CI.] Un pipeline Python
    qui interroge les API Azure, Wazuh et l'AD, normalise les
    réponses, détecte les changements, et régénère le dossier
    d'homologation au format ANSSI. À chaque exécution, on a un
    nouveau dossier daté, signé, traçable.
\end{description}

\section{Ce qu'on cherche à vérifier}

Trois objectifs concrets pour l'outil~:

\begin{enumerate}
  \item \textbf{Couverture}~: quelle proportion des preuves
    techniques attendues dans un dossier d'homologation ANSSI
    l'outil arrive-t-il à collecter tout seul~?
  \item \textbf{Fraîcheur}~: combien de temps s'écoule entre un
    changement réel (par exemple, une règle NSG modifiée) et la
    mise à jour du dossier~?
  \item \textbf{Précision}~: ce que l'outil écrit dans le
    dossier correspond-il vraiment à ce qui tourne dans le lab~?
\end{enumerate}

Ce ne sont pas des hypothèses scientifiques au sens strict, ce sont
des objectifs d'ingénieur. On les mesure dans la partie~IV.

\section{Ce que ce mémoire n'est pas}

Pour cadrer dès le départ~:

\begin{itemize}
  \item \textbf{Ce n'est pas une refonte de la méthode ANSSI.}
    Le guide en neuf étapes reste la référence. L'outil produit
    un dossier qui suit cette structure, pas un format inventé.
  \item \textbf{Ce n'est pas un produit commercial.} C'est un PoC
    fonctionnel, codé proprement, mais sans support, sans
    interface graphique, sans intégration au-delà du lab.
  \item \textbf{Ce n'est pas une couverture exhaustive.} Tout ce
    qui est organisationnel (formation des équipes, procédures,
    revues humaines) reste manuel. L'outil n'attaque que la
    couche technique.
\end{itemize}

\section{Plan du mémoire}

Le mémoire suit cinq parties~:

\begin{description}
  \item[Partie~I --- Contexte.] Le problème métier, et un tour de
    l'existant côté audit et conformité continue
    (\cref{ch:etat-existant}).
  \item[Partie~II --- Le lab.] Comment UrbanUpC est monté~:
    architecture Azure (\cref{ch:lab-archi}), applications et AD
    (\cref{ch:apps-ad}), SOC Wazuh (\cref{ch:soc-wazuh}).
  \item[Partie~III --- L'outil.] Ce que HOMO-CI doit faire
    (\cref{ch:conception}) et comment il est codé
    (\cref{ch:implementation}).
  \item[Partie~IV --- Ce qu'on a mesuré.] Les tests sur le lab et
    les chiffres réels (\cref{ch:resultats}).
  \item[Partie~V --- Discussion et suite.] Les limites honnêtes,
    ce qui marcherait au-delà du lab, et la conclusion
    (\cref{ch:limites}).
\end{description}
